% -*- coding: utf-8 -*-
% !TEX program = xelatex
\documentclass[plain,noanswer,medmath]{../../template/jnuexam}
\usepackage{../../template/kysx}

\begin{document}


\examtitle{2010}{1}


\makepart{选择题}{1～8小题，每小题4分，共32分}


\begin{problem}
极限$\lim_{x \to \infty} \left[\frac{x^2}{(x - a)(x + b)} \right]^x =$\pickout{}
\begin{abcd}
\item $1$．
\item $\e$．
\item $\e^{a - b}$．
\item $\e^{b - a}$．
\end{abcd}
\end{problem}


\begin{problem}
设函数$z = z(x,y)$由方程$F\left(\frac{y}{x},\frac{z}{x} \right) = 0$确定，其中$F$为可微函数，
且$F'_2 \ne 0$，则$x\frac{\pdz}{\pdx} + y\frac{\pdz}{\pdy} =$\pickout{}
\begin{abcd}
\item $x$．
\item $z$．
\item $- x$．
\item $- z$．
\end{abcd}
\end{problem}


\begin{problem}
设$m,n$是正整数，则反常积分$\int_0^1 \frac{\sqrt[m]{\ln^2(1 - x)}}{\sqrt[n]{x}}\dx$的收敛性\pickout{}
\begin{abcd}
\item 仅与$m$的取值有关．
\item 仅与$n$的取值有关．
\item 与$m,n$取值都有关．
\item 与$m,n$取值都无关．
\end{abcd}
\end{problem}


\begin{problem}
$\lim_{n\to\infty} \sum_{i=1}^n \sum_{j=1}^n \frac{n}{(n+i)\left(n^2 + j^2 \right)} =$\pickout{}
\begin{abcd}
\item $\int_0^1 \dx\int_0^x \frac{1}{(1 + x)\left(1 + y^2 \right)} \dy$．
\item $\int_0^1 \dx\int_0^x \frac{1}{(1 + x)\left(1 + y \right)} \dy$．
\item $\int_0^1 \dx\int_0^1 \frac{1}{(1 + x)\left(1 + y \right)} \dy$．
\item $\int_0^1 \dx\int_0^1 \frac{1}{(1 + x)\left(1 + y^2 \right)} \dy$．
\end{abcd}
\end{problem}


\begin{problem}
设$A$为$m \times n$矩阵，$B$为$n \times m$矩阵，$E$为$m$阶单位矩阵，若$AB = E$，则\pickout{}
\begin{abcd}
\item 秩$r(A) = m$，秩$r(B) = m$．
\item 秩$r(A) = m$，秩$r(B) = n$．
\item 秩$r(A) = n$，秩$r(B) = m$．
\item 秩$r(A) = n$，秩$r(B) = n$．
\end{abcd}
\end{problem}


\begin{problem}
设$A$为$4$阶实对称矩阵，且$A^2 + A = O$，若$A$的秩为$3$，则$A$相似于\pickout{}
\begin{abcd}
\item $\begin{pmatrix}
  1 &   &   & \\
    & 1 &   & \\
    &   & 1 & \\
    &   &   & 0
\end{pmatrix}$．
\item $\begin{pmatrix}
  1 &   &    & \\
    & 1 &    & \\
    &   & -1 & \\
    &   &    & 0
\end{pmatrix}$．
\item $\begin{pmatrix}
  1 &    &    & \\
    & -1 &    & \\
    &    & -1 & \\
    &    &    & 0
\end{pmatrix}$．
\item $\begin{pmatrix}
  -1 &    &    & \\
     & -1 &    & \\
     &    & -1 & \\
     &    &    & 0
\end{pmatrix}$．
\end{abcd}
\end{problem}


\begin{problem}
设随机变量$X$的分布函数$F(x) = \begin{cases}
             0, & x < 0 \\
   \frac{1}{2}, & 0 \le x < 1 \\
  1 - \e^{- x}, & x \ge 1
\end{cases}$，则$P\{X = 1\} =$\pickout{}
\begin{abcd}
\item $0$．
\item $\frac{1}{2}$．
\item $\frac{1}{2} - \e^{-1}$．
\item $1 - \e^{-1}$．
\end{abcd}
\end{problem}


\begin{problem}
设$f_1(x)$为标准正态分布的概率密度，$f_2(x)$为$\left[- 1,3 \right]$上均匀分布的概率密度，若$f(x) = \begin{cases}
  af_1(x), & x \le 0 \\
  bf_2(x), & x > 0
\end{cases}$$(a > 0,b > 0)$为概率密度，则$a,b$应满足\pickout{}
\begin{abcd}
\item $2a + 3b = 4$．
\item $3a + 2b = 4$．
\item $a + b = 1$．
\item $a + b = 2$．
\end{abcd}
\end{problem}


\makepart{填空题}{9～14小题，每小题4分，共24分}


\begin{problem}
设$\begin{cases}
  x = \e^{- t}, \\
  y = \int_0^t \ln \left(1 + u^2 \right)\du,
\end{cases}$求$\left. \frac{\d^2y}{\dx^2} \right|_{t = 0} =$ \fillin{}．
\end{problem}


\begin{problem}
$\int_0^{\pi^2} \sqrt{x} \cos \sqrt{x} \dx =$ \fillin{}．
\end{problem}


\begin{problem}
已知曲线$L$的方程为$y = 1 - |x|$（$x\in[-1,1]$），起点是$(- 1,0)$，终点是$(1,0)$，
则曲线积分$\int_L xy\dx + x^2 \dy =$ \fillin{}．
\end{problem}


\begin{problem}
设$\Omega = \left\{\left(x,y,z \right)\left| x^2 + y^2 \le z \le 1 \right. \right\}$，
则$\Omega$的形心的竖坐标$\overline{z} =$ \fillin{}．
\end{problem}


\begin{problem}
设$\alpha_1 = \left(1,2, - 1,0 \right)^T$, $\alpha_2 = \left(1,1,0,2 \right)^T$,
$\alpha_3 = \left(2,1,1,a \right)^T$，若由$\alpha_1,\alpha_2,\alpha_3$生成的向量空间的维数是$2$，
则$a =$ \fillin{}．
\end{problem}


\begin{problem}
设随机变量$X$的概率分布为$P\{X=k\} = \frac{C}{k!}$, $k=0,1,2,\cdots$，则$E\left(X^2\right)=$ \fillin{}．
\end{problem}


\makepart{解答题}{15～23小题，共94分}


\begin{problem}[points=10]
求微分方程$y'' - 3y' + 2y = 2x\e^x$的通解．
\end{problem}


\begin{problem}[points=10]
求函数$f(x) = \int_1^{x^2} \left(x^2 - t \right) \e^{- t^2} \dt$的单调区间与极值．
\end{problem}


\begin{problem}[points=10]
\begin{enumerate}
\item
比较$\int_0^1 |\ln t|[\ln(1+t)]^n\dt$
与$\int_0^1 t^n |\ln t|\dt$（$n=1,2,\cdots$）的大小，说明理由；
\item
记$u_n = \int_0^1 |\ln t|[\ln(1+t)]^n\dt$（$n=1,2,\cdots$），求极限$\lim_{n\to\infty} u_n$．
\end{enumerate}
\end{problem}


\begin{problem}[points=10]
求幂级数$\sum_{n = 1}^{\infty}\frac{(- 1)^{n - 1}}{2n - 1}x^{2n}$的收敛域及和函数．
\end{problem}


\begin{problem}[points=10]
设$P$为椭球面$S: x^2 + y^2 + z^2 - yz = 1$上的动点，若$S$在点$P$处的切平面与$xOy$面垂直，
求点$P$的轨迹$C$，并计算曲面积分
$I = \iint_{\Sigma}\frac{\left(x + \sqrt{3} \right)|y - 2z|}{\sqrt{4 + y^2 + z^2 - 4yz}} \dS$，
其中$\Sigma$是椭球面$S$位于曲线$C$上方的部分．
\end{problem}


\begin{problem}[points=11]
设$A = \begin{pmatrix}
  \lambda & 1           & 1 \\
  0       & \lambda - 1 & 0 \\
  1       & 1           & \lambda
\end{pmatrix}$, $b = \begin{pmatrix}
  a \\
  1 \\
  1
\end{pmatrix}$，已知线性方程组$Ax = b$存在两个不同的解．\par
\begin{enumerate*}
\item 求$\lambda$, $a$；
\item 求方程组$Ax = b$的通解．
\end{enumerate*}
\end{problem}


\begin{problem}[points=11]
已知二次型$f(x_1,x_2,x_3) = x^T Ax$在正交变换$x = Qy$下的标准形为$y_1^2 + y_2^2$，
且$Q$的第$3$列为$\left(\frac{\sqrt{2}}{2},0,\frac{\sqrt{2}}{2} \right)^T$．\par
\begin{enumerate*}
\item 求矩阵$A$；
\item 证明$A + E$为正定矩阵，其中$E$为$3$阶单位矩阵．
\end{enumerate*}
\end{problem}


\begin{problem}[points=11]
设二维随机变量$(X,Y)$的概率密度为
$$f(x,y) = A\e^{- 2x^2 + 2xy - y^2},\quad -\infty < x < +\infty , -\infty < y < +\infty .$$
求常数$A$及条件概率密度$f_{Y|X}(y|x)$．
\end{problem}


\begin{problem}[points=11]
设总体$X$的概率分布为
$$\begin{array}{|c|ccc|}
\hline
  X &          1 &                 2 & 3 \\
\hline
  P & 1 - \theta & \theta - \theta^2 & \theta^2 \\
\hline
\end{array}$$
其中参数$\theta \in \left(0,1 \right)$未知，
以$N_i$表示来自总体$X$的简单随机样本（样本容量为$n$）中等于$i$的个数（$i = 1,2,3$）．
试求常数$a_1$, $a_2$, $a_3$，使$T = \sum_{i = 1}^3 a_i N_i$为$\theta$的无偏估计量，并求$T$的方差．
\end{problem}


\end{document}
